%!TEX encoding = UTF8
%!TEX root = 0-notes.tex

\chapter*{Préface}


\section*{Syllabus}

\subsection*{Description}

\subsection*{Objectifs}

\subsection*{Pré-requis nécessaires}

\subsection*{Pré-requis recommandés}

\section*{Structure du cours}

Construction d'un cours de mathématiques : Exemple -- Remarque -- Définition -- Lemme -- Proposition -- Théorème -- Démonstration -- Corollaire.

Voir \emph{Mechanica, volume 1}, Leonhard Euler, 1736 (\href{https://scholarlycommons.pacific.edu/euler-works/15/}{E15}), pour un exemple de cours de mathématiques en latin.

Wolframalpha permet de faire du calcul exact.

\begin{definition*}
Associe un mot de vocabulaire à un objet mathématique.
Un mot de vocabulaire peut avoir un sens différent en mathématiques qu'en langage vernaculaire.
C'est par exemple le cas de « ensemble », « couple », « rationnel », « réel », ...
\end{definition*}

\begin{lemme*}
	Résultat préliminaire, parfois technique, utile à la démonstration d'une proposition plus importante.
\end{lemme*}

\begin{proposition*}
	Énoncé vrai et important.
\end{proposition*}

\begin{theorem*}
	Énoncé vrai et fondamental.
\end{theorem*}

\begin{corollaire*}
	Résultat particulier, découlant souvent quasi directement d'un théorème ou d'une proposition.
\end{corollaire*}

%\newcounter{preface}
\begin{Exercise}[label=1]%[counter=preface]
	Exemple d'exercice d'application directe.
\end{Exercise}
\begin{Exercise}[difficulty=1, label=2]%, counter=preface]
	Exemple de problème.
\end{Exercise}
\begin{Exercise}[difficulty=2, label=3]%, counter=preface]
	Exemple de problème d'approfondissement.
\end{Exercise}
\shipoutExercise
\begin{Answer}[ref=1]
	Solution 1.
\end{Answer}
\begin{Answer}[ref=2]
	Solution 2.
\end{Answer}
\begin{Answer}[ref=3]
	Solution 3.
\end{Answer}
\begin{multicols}{3}
\shipoutAnswer
\end{multicols}
